\documentclass[sigconf,nonacm]{acmart}

\usepackage{booktabs}
\usepackage{longtable}
\usepackage{listings}
\usepackage{xcolor}
\usepackage{hyperref}
\usepackage{enumitem}
\usepackage{caption}
\usepackage{subcaption}

\definecolor{codebg}{RGB}{245,245,245}
\definecolor{codeframe}{RGB}{200,200,200}

\lstset{
    basicstyle=\small\ttfamily,
    breaklines=true,
    breakatwhitespace=true,
    columns=flexible,
    keepspaces=true,
    showstringspaces=false,
    frame=single,
    rulecolor=\color{codeframe},
    backgroundcolor=\color{codebg},
    xleftmargin=1em,
    xrightmargin=1em,
    aboveskip=0.5em,
    belowskip=0.5em,
}

\captionsetup[table]{skip=4pt}
\captionsetup[figure]{skip=4pt}

\begin{document}

\title{VeritasCarbon: Appendix Materials}
\subtitle{Supplementary Documentation for VLDB 2027 Submission}

\maketitle

\section*{Dataset as a Core Contribution}

In addition to the CoDE multi-agent framework, this paper makes a substantial dataset contribution by releasing \textbf{VeritasCarbon-ESG-35K}, the largest publicly available traceable ESG instruction dataset to date. The dataset contains 35,009 high-quality instruction--response pairs generated from 17,721 real-world ESG disclosure documents spanning four corpus layers (domain textbooks, CSR reports, regulatory guidelines, and industry analyses). Every pair is fully traceable to its source document chunk and includes rich generation metadata (quality score, selected experts, collaboration mode, knowledge injection count), enabling fine-grained analysis of data provenance and quality.

Unlike existing ESG corpora that are limited to classification labels or short extracts, VeritasCarbon-ESG-35K provides \textit{task-oriented instructions} with diverse types (Q\&A, summarization, extraction, classification, analysis, temporal analysis, benchmark comparison, greenwashing detection, standard alignment, knowledge graph construction, and consistency verification). This makes it directly usable for instruction tuning of large language models in the ESG and sustainable-finance domains.

The dataset is released under a tiered strategy to balance accessibility and copyright constraints:
\begin{itemize}[noitemsep]
    \item A representative 2,000-pair sample is included in the GitHub repository for immediate inspection and replication of the main evaluation (Table~2);
    \item The full 35,009-pair dataset is hosted on Hugging Face Datasets (\url{https://huggingface.co/datasets/Yihan-JIANG/VeritasCarbon-ESG-35K}) under CC BY-SA~4.0;
    \item Complete provenance metadata for all 17,721 source documents is provided in the repository, and the full dataset can be reproduced from source using the provided notebooks.
\end{itemize}

\appendix

\section{Data and Code Availability}

\noindent\textbf{Code.} The complete VeritasCarbon framework, evaluation scripts, and reproducibility instructions are available at \url{https://github.com/YihanJIANG-lab/VeritasCarbon-VLDB2027} under the MIT License.

\noindent\textbf{Sample Data.} A representative 2,000-pair sample (\texttt{veritascarbon\_sample\_2000.jsonl}, 8.7~MB) is included in the GitHub repository. It is drawn with \texttt{random.seed(42)} from the full 35,009-pair pool and matches the evaluation protocol used in Table~2.

\noindent\textbf{Full Dataset.} The complete VeritasCarbon-ESG-35K dataset (35,009 instruction--response pairs, $\sim$153~MB) is hosted on Hugging Face Datasets:
\begin{center}
\url{https://huggingface.co/datasets/Yihan-JIANG/VeritasCarbon-ESG-35K}
\end{center}
It can be loaded in one line:
\begin{lstlisting}[language=Python]
from datasets import load_dataset
ds = load_dataset("Yihan-JIANG/VeritasCarbon-ESG-35K", split="train")
\end{lstlisting}

\noindent\textbf{Raw Corpus.} The source ESG documents ($\sim$2.7~GB, 17,721 files) contain copyrighted material and cannot be redistributed in full. Complete provenance metadata (document IDs, titles, source URLs, and chunk-to-document mappings) is provided in \texttt{data/raw\_corpus/CORPUS\_MANIFEST.md}.

\section{Reproducibility Checklist}

Table~\ref{tab:reproducibility} provides a 30-minute verification checklist for reviewers. All commands run on CPU; no GPU is required for verification.

\begin{table}[h]
\centering
\caption{30-Minute Reproducibility Checklist}
\label{tab:reproducibility}
\small
\begin{tabular}{@{}p{3.5cm}p{5.2cm}p{2.5cm}r@{}}
\toprule
\textbf{Claim} & \textbf{Verification Command} & \textbf{Expected Output} & \textbf{Time} \\
\midrule
Table~2: CoDE ROUGE-L & \texttt{cat results/outputs/intrinsic\_comparison.json} & 0.3380 & 30~s \\
Table~2: CoDE BLEU-4 & same file & 0.1932 & --- \\
Table~2: Domain Rel. & same file & 0.3637 & --- \\
Table~2: FactCheck & same file & 0.9478 & --- \\
Table~3: Ablation & \texttt{cat results/figures\_and\_tables/table3\_ablation.tex} & all 12 values & 30~s \\
Fig~5: Quality dist. & \texttt{python scripts/generate\_paper\_results.py --figures} & PDF/PNG files & 2~min \\
Dataset: 35,009 & \texttt{cat results/figures\_and\_tables/table2\_qa\_statistics.json} & 35009 & 10~s \\
Provenance & Inspect sample \texttt{jsonl} & \texttt{source\_chunk\_id} present & 1~min \\
Eval pipeline & Run \texttt{intrinsic\_evaluation\_03\_03.py} & metrics JSON & 10~min \\
\bottomrule
\end{tabular}
\end{table}

\noindent\textbf{Notes.} (1)~Tracks A--B require only the GitHub repository. (2)~Track~C requires the raw corpus and an A100 GPU; see \texttt{REPRODUCIBILITY.md} for details. (3)~All random sampling uses \texttt{random.seed(42)}.

\section{Prompt Templates}

The CoDE framework employs 11 specialized expert agents plus a MetaExpert orchestrator. Each expert receives a system prompt that defines its persona, task, and output format. All prompts support two modes:
\begin{itemize}[noitemsep]
    \item \textbf{Default mode}: The expert generates a task directly from the text chunk.
    \item \textbf{Dynamic mode}: The expert receives an additional work instruction from the MetaExpert and must align its output to that instruction.
\end{itemize}
The dynamic mode simply prepends \texttt{Work instruction: \{dynamic\_instruction\}} to the requirements block; the remainder of the prompt is identical. Below we present the default-mode prompts for representative experts, followed by a summary table for all 11 experts, and finally the MetaExpert orchestration prompt.

\subsection{Q\&A Expert}
\begin{lstlisting}
You are a senior ESG (environment, social, governance) Q&A generation expert.
Task: Generate one high-quality Q&A pair from the given text.
Requirements:
1. The question should be highly relevant to the text.
2. The answer must be findable in the text or reasonably inferred.
3. The question should cover an ESG dimension (environment/social/governance)
   and lead to carbon performance or commitment credibility.
4. The question should be specific and clear.
5. The answer should be accurate and complete.
{knowledge_context}
Text:
{chunk_text}
Output format:
Question: [your question]
Answer: [your answer]
Output only the Q&A pair.
\end{lstlisting}

\subsection{Analysis Expert}
\begin{lstlisting}
You are a senior ESG analysis expert.
Task: Generate one analysis task (instruction-answer pair) from the text.
Requirements:
1. The instruction should ask to analyze ESG content (e.g. comparison,
   trend, impact) and lead to carbon/commitment assessment.
2. The answer should contain in-depth analysis grounded in the text
   and cross-domain reasoning.
3. Be critical and assess credibility of carbon commitments.
{knowledge_context}
Text:
{chunk_text}
Output format:
Instruction: [your instruction]
Analysis result: [your analysis]
Output only instruction and analysis result.
\end{lstlisting}

\subsection{Greenwashing Detection Expert}
\begin{lstlisting}
You are a senior ESG greenwashing detection expert, skilled at identifying
misleading or exaggerated claims in reports.
Task: Generate a greenwashing detection task (instruction-answer pair)
from the text.
Requirements:
1. Identify greenwashing patterns (refer to common patterns list).
2. Quantify greenwashing risk (0-1, 1 = highest risk).
3. Provide evidence chain (quote source text).
4. Produce greenwashing risk report.

Common greenwashing patterns:
1. Vague commitments: "committed to", "plan to" without concrete targets
2. Selective disclosure: only favorable data, hiding unfavorable info
3. Exaggerated impact: overstating environmental project effects
4. Time mismatch: presenting future plans as achieved results
5. Scope confusion: mixing scope 1/2/3 emissions
6. Inconsistency: commitments vs actual actions
7. Lack of verification: no third-party or traceable evidence
{knowledge_context}
Text:
{chunk_text}
Output format:
Instruction: [detection instruction]
Answer: [matched patterns, risk level, evidence chain]
Output only instruction and answer.
\end{lstlisting}

\subsection{Consistency Verification Expert}
\begin{lstlisting}
You are a senior ESG consistency verification expert, skilled at detecting
data contradictions and gaps in reports.
Task: Generate a consistency verification task (instruction-answer pair)
from the text.
Requirements:
1. Check whether the same data is stated consistently across the text.
2. Identify contradictions (e.g. scope-1 emissions differing across sections).
3. Detect "missing data" patterns (only favorable data disclosed).
4. Produce consistency report.
{knowledge_context}
Text:
{chunk_text}
Output format:
Instruction: [verification instruction]
Answer: [consistency assessment, contradictions, missing patterns]
Output only instruction and answer.
\end{lstlisting}

\subsection{Summary of All 11 Experts}

Table~\ref{tab:experts} summarizes the role and output format of each expert agent. All experts share the same prompt structure (persona + task + requirements + knowledge context + chunk text + output format); only the specific task description and output labels differ.

\begin{table}[h]
\centering
\caption{CoDE Expert Agent Summary}
\label{tab:experts}
\small
\begin{tabular}{@{}p{3cm}p{5.5cm}p{2.2cm}@{}}
\toprule
\textbf{Expert} & \textbf{Core Task} & \textbf{Output Labels} \\
\midrule
Q\&A Expert & Generate a specific Q\&A pair grounded in the text & Question / Answer \\
Summary Expert & Summarization (100--200 words) highlighting ESG & Instr / Summary \\
Extraction Expert & Extract structured ESG information & Instr / Extraction \\
Classification Expert & Classify by ESG dimension (E/S/G/Mixed) & Instr / Classification \\
Analysis Expert & In-depth cross-domain critical assessment & Instr / Analysis \\
Temporal Analysis Expert & Multi-year trends and commitment consistency & Instr / Answer \\
Benchmark Comparison Expert & Compare against industry averages & Instr / Answer \\
Greenwashing Detection Expert & Identify misleading claims and quantify risk & Instr / Answer \\
Standard Alignment Expert & Identify standards (GRI, TCFD, SASB) and gaps & Instr / Answer \\
Knowledge Graph Expert & Extract entities/relations; detect broken links & Instr / Answer \\
Consistency Verification Expert & Detect contradictions and selective disclosure & Instr / Answer \\
\bottomrule
\end{tabular}
\end{table}

\subsection{MetaExpert Orchestrator}

The MetaExpert designs work instructions for subordinate experts, evaluates their outputs, and provides iterative feedback. Its prompt dynamically incorporates the target expert type, core topics, chunk preview, and collaboration context.

\begin{lstlisting}
# Role
You are a top carbon-strategy analyst. Your task is to evaluate all
information related to the company's carbon emissions and carbon-reduction
commitments. Now design one work instruction for a subordinate analyst.

# Context
- Subordinate role: {expert_type} ({expert_desc})
- Chunk is from report section [{section_name}]
- Core topics: {topics_str}
- Preview: {chunk_text[:500]}...
- Collaboration: This run uses multiple experts: {all_experts}
  Current expert: {expert_type}; Others: {other_experts}
  Requirement: Your instruction should match {expert_type}'s role and
  complement others (avoid duplication).

# Task
Generate one work instruction for this analyst. The instruction must:
1. Carbon-centric: lead to assessing carbon performance or commitment credibility.
2. Cross-domain: use information in this chunk as evidence for carbon strategy.
3. Depth: guide the analyst to compare, evaluate, predict, or critique
   -- not just extract.

# Output
Work instruction:
\end{lstlisting}

\noindent\textbf{Feedback round.} After receiving the expert's output, the MetaExpert runs a quality-assessment prompt that scores correctness, completeness, relevance, and structure on a 0--1 scale. If any dimension falls below threshold, a refined work instruction is generated and the expert re-executes, up to $R=2$ rounds.

\section{Hyperparameters}

Table~\ref{tab:hyperparameters} lists all fixed hyperparameters used in the experiments. All random sampling steps use \texttt{random.seed(42)}.

\begin{table}[h]
\centering
\caption{Complete Hyperparameter Configuration}
\label{tab:hyperparameters}
\small
\begin{tabular}{@{}p{4cm}p{3cm}p{4.5cm}@{}}
\toprule
\textbf{Parameter} & \textbf{Value} & \textbf{Source File} \\
\midrule
Base model & Qwen2-72B-Instruct & \texttt{configs/config.yaml} \\
Quantization & 4-bit (NF4) & Unsloth default \\
Temperature & 0.7 & \texttt{coe\_framework\_02\_03.py} \\
Max tokens (instruction) & 512 & --- \\
Max tokens (response) & 1,024 & --- \\
Chunk size & 2,000 chars & \texttt{text\_chunker\_01\_03.py} \\
Chunk overlap & 200 chars & \texttt{text\_chunker\_01\_03.py} \\
Expert selection threshold & 0.5 & \texttt{expert\_selector\_02\_01.py} \\
Collaboration mode & parallel / sequential & \texttt{coe\_framework\_02\_03.py} \\
Feedback rounds ($R$) & 2 (default) & \texttt{meta\_expert\_02\_09.py} \\
Quality threshold (accept) & 0.6 & \texttt{meta\_expert\_02\_09.py} \\
Knowledge items per chunk & up to 5 & \texttt{domain\_knowledge\_02\_02.py} \\
Evaluation sample size & 2,000 pairs & \texttt{intrinsic\_evaluation\_03\_03.py} \\
Random seed & 42 & All sampling steps \\
\bottomrule
\end{tabular}
\end{table}

\section{Expert Selection Logic}

Expert selection is performed by a rule-based feature extractor (\texttt{expert\_selector\_02\_01.py}) that scores each chunk against predefined feature patterns. No LLM is invoked during selection; the process is fully deterministic.

\begin{enumerate}[noitemsep]
    \item \textbf{Feature extraction.} The chunk is analyzed for:
    \begin{itemize}[noitemsep]
        \item Keyword presence (carbon, emission, energy, employee, training, governance, etc.)
        \item Numerical data density (counts, percentages, monetary values)
        \item Temporal markers (years, quarters, ``since'', ``by 20xx'')
        \item Standard references (GRI, TCFD, SASB, ISO 14001, CDP)
        \item Comparative language (``higher than'', ``below average'', ``ranked'')
        \item Entity types (company names, plant locations, regulatory bodies)
    \end{itemize}
    
    \item \textbf{Expert scoring.} Each expert defines $s_i(\text{chunk}) \in [0,1]$. Examples:
    \begin{itemize}[noitemsep]
        \item \textit{TemporalAnalysisExpert}: high score if temporal markers $\geq 2$ and numerical trends present.
        \item \textit{GreenwashingDetectionExpert}: high score if vague-commitment keywords + selective-disclosure patterns detected.
        \item \textit{KnowledgeGraphExpert}: high score if entity count $\geq 5$ and relation verbs present.
    \end{itemize}
    
    \item \textbf{Thresholding.} Experts with $s_i \geq 0.5$ are selected; fallback is QAExpert.
    
    \item \textbf{Collaboration mode.} If $|\{s_i \geq 0.5\}| \geq 2$, parallel; otherwise sequential.
\end{enumerate}

\noindent\textbf{MetaExpert override.} After selection, the MetaExpert may dynamically adjust the assignment based on core topics. This override is rare ($<5\%$) and logged in metadata.

\section{Error Analysis}

To illustrate the quality spectrum, we present representative examples from the 2,000-pair sample. The quality score is a composite metric (correctness 0.4, completeness 0.3, relevance 0.2, structure 0.1) assigned by the MetaExpert.

\subsection{High-Quality Example (Score 0.915)}

\textbf{Experts:} Summary + Consistency Verification\\
\textbf{Instruction:} Generate a summary about 2024--2025 A-share ESG fund issuance trends, disclosure quality, carbon-risk considerations, emission management progress, evaluation-system effectiveness, and green-transformation potential.

\begin{lstlisting}
Answer: Since 2024, A-share ESG public fund issuance has reached
approximately 216 billion yuan. Although the issuance pace has slowed,
18 new funds were still issued from Q4 to Q3. The industry scale has
rebounded against the trend, showing a structural shift from quantity
expansion to quality improvement. Under policy and market dual-wheel
drive, China's ESG market is gradually maturing.
\end{lstlisting}

\noindent\textbf{Why high:} The instruction demands multi-faceted synthesis, and the answer integrates quantitative evidence (216B yuan, 18 funds) with qualitative assessment. Consistency Verification cross-checks fund-scale figures.

\subsection{Low-Quality Example (Score 0.000)}

\textbf{Experts:} Summary\\
\textbf{Instruction:} Generate a summary for the following text.

\begin{lstlisting}
Answer: This article aims to analyze how ESG information disclosure
affects employees' green innovation efficiency and its role in promoting
companies' real carbon-reduction commitments. Special attention is paid
to identifying and criticizing potential "greenwashing" behavior...
\end{lstlisting}

\noindent\textbf{Why low:} The output is a generic abstract rather than a task-oriented pair. The instruction is merely ``generate a summary,'' and the answer restates the paper's purpose without concrete data. This occurs when the MetaExpert's work instruction is too vague and no second expert refines depth.

\subsection{Key Failure Patterns}

Table~\ref{tab:failure} summarizes common failure modes in the bottom-decile records.

\begin{table}[h]
\centering
\caption{Common Failure Patterns in Low-Quality Records}
\label{tab:failure}
\small
\begin{tabular}{@{}p{2.5cm}p{4.5cm}p{3.5cm}@{}}
\toprule
\textbf{Pattern} & \textbf{Description} & \textbf{Typical Cause} \\
\midrule
Generic summary & Answer restates source abstract without specificity & Single Summary Expert, no dynamic instruction \\
Vague instruction & Broad request (``analyze ESG'') & MetaExpert topics too generic \\
Missing evidence & Claims lack concrete numbers & No Extraction/Benchmark Expert \\
Self-reference & Answer refers to ``this article'' & Expert fails to rephrase \\
\bottomrule
\end{tabular}
\end{table}

\section{Dataset Schema}

Each record is a JSON object with the schema in Table~\ref{tab:schema}.

\begin{table}[h]
\centering
\caption{VeritasCarbon-ESG-35K JSON Schema}
\label{tab:schema}
\small
\begin{tabular}{@{}p{2.5cm}p{1.5cm}p{6.5cm}@{}}
\toprule
\textbf{Field} & \textbf{Type} & \textbf{Description} \\
\midrule
\texttt{instruction} & string & The generated task instruction (avg 106.5 chars) \\
\texttt{input} & string & Source chunk text (up to $\sim$2,000 chars) \\
\texttt{output} & string & Generated response (avg 380.4 chars) \\
\texttt{source\_chunk\_id} & string & Unique chunk identifier \\
\texttt{metadata} & object & Generation provenance and quality metadata \\
\bottomrule
\end{tabular}
\end{table}

\subsection{Metadata Object}

\begin{table}[h]
\centering
\small
\begin{tabular}{@{}p{3.2cm}p{1.5cm}p{5.8cm}@{}}
\toprule
\textbf{Field} & \textbf{Type} & \textbf{Description} \\
\midrule
\texttt{quality\_score} & float & Composite score $[0,1]$ (mean 0.667) \\
\texttt{selected\_experts} & list & IDs of invoked agents \\
\texttt{expert\_names} & string & Human-readable names \\
\texttt{selection\_reasons} & object & Feature scores triggering selection \\
\texttt{use\_meta\_expert} & bool & Whether MetaExpert was used \\
\texttt{dynamic\_instructions} & string & MetaExpert-generated instruction \\
\texttt{knowledge\_items\_count} & int & Injected knowledge items (0--5) \\
\texttt{collaboration\_mode} & string & \texttt{parallel} or \texttt{sequential} \\
\texttt{max\_iterations} & int & Feedback rounds (0--2) \\
\texttt{model} & string & Qwen2-72B-Instruct \\
\texttt{framework} & string & CoDE \\
\bottomrule
\end{tabular}
\end{table}

\noindent\textbf{Quality composition:} Correctness (0.4) + Completeness (0.3) + Relevance (0.2) + Structure (0.1). Score $<0.6$ triggers feedback; still failing after $R=2$ rounds leads to discard.

\section{Latency and Cost Breakdown}

Table~\ref{tab:latency} reports end-to-end latency on a single NVIDIA A100 (80~GB) with Qwen2-72B-Instruct 4-bit quantized via Unsloth.

\begin{table}[h]
\centering
\caption{Latency and Cost Breakdown}
\label{tab:latency}
\small
\begin{tabular}{@{}lrrr@{}}
\toprule
\textbf{Stage} & \textbf{Time / Rec} & \textbf{Total (35K)} & \textbf{\%} \\
\midrule
Document parsing & --- & $\sim$2~hrs & 4\% \\
Semantic chunking & --- & $\sim$30~min & 1\% \\
Expert selection & 0.02~s & $\sim$12~min & $<$1\% \\
Knowledge retrieval & 0.15~s & $\sim$88~min & 3\% \\
Parallel generation & 2.8~s & $\sim$27~hrs & 58\% \\
MetaExpert feedback ($R=2$) & 1.9~s & $\sim$18~hrs & 38\% \\
Quality scoring & 0.05~s & $\sim$29~min & 1\% \\
\midrule
\textbf{Total} & \textbf{$\sim$4.9~s} & \textbf{$\sim$46~hrs} & \textbf{100\%} \\
\bottomrule
\end{tabular}
\end{table}

\noindent\textbf{Notes.} (1)~Parsing and chunking are CPU-bound. (2)~Generation and feedback are GPU-bound and dominate. (3)~With 4 A100 GPUs in parallel, wall-clock time drops to $\sim$12 hours.

\end{document}
